\subsection{线性加权模型的关系}
\label{sec:regression-linear-relation}

本节各模型共享“对特征作线性加权”的骨架，差别只在四个可独立变化的组件：残差怎样计分、参数怎样受约束、特征怎样表示、响应结构怎样指定。表\ref{tab:regression-linear-comparison}按这四个组件排列。

\begin{table*}[t]
  \centering
  \small
  \setlength{\tabcolsep}{4pt}
  \renewcommand{\arraystretch}{1.25}
  \begin{tabularx}{\linewidth}{@{}>{\RaggedRight\arraybackslash}p{22mm}>{\RaggedRight\arraybackslash}p{30mm}>{\RaggedRight\arraybackslash}p{26mm}>{\RaggedRight\arraybackslash}X@{}}
    \toprule
    模型 & 残差计分 & 参数约束 & 特征与响应结构 \\
    \midrule
    OLS & 平方 & 无 & 原始输入，高斯同方差 \\
    \addlinespace
    GLS & 加权平方 & 无 & 已知误差协方差加权 \\
    \addlinespace
    岭/Lasso/弹性网/组Lasso & 平方 & 二次/一范数/混合/组范数 & 原始输入，可加惩罚 \\
    \addlinespace
    PCR/PLS & 平方 & 成分截断 & 监督或无监督降维后拟合 \\
    \addlinespace
    Huber & 分段平方+线性 & 无 & 稳健计分，不指定分布 \\
    \addlinespace
    分位数 & 非对称加权 & 无（可加一范数） & 估计条件分位，非高斯 \\
    \addlinespace
    SVR & $\epsilon$不敏感 & 二次范数 & 可换核内积 \\
    \addlinespace
    基函数/MARS/KRR & 平方 & 基内或核空间惩罚 & 固定基、数据选基或隐式核 \\
    \addlinespace
    GLM/GAM & 指数族对数似然 & 可加惩罚/平滑 & 指数族+链接，可加平滑项 \\
    \bottomrule
  \end{tabularx}
  \caption{线性加权模型的组件比较。同一行内各模型的损失、惩罚与结构假设不能互相替代；表中只列主要组件，超参数与验证流程仍按各模型规定。}
  \label{tab:regression-linear-comparison}
\end{table*}

用相同输入表示比较OLS、Huber、分位数与SVR的残差计分和预测含义：异常响应由Huber的分段线性降权，非对称代价由分位数的非对称加权表达，容忍误差由SVR的$\epsilon$带吸收，三者回答不同决策问题。

再比较基函数、成分方向、系数惩罚、精度加权和响应链接的作用。PCR按输入方差选方向、PLS让响应参与逐步构造方向，多项式与样条通过增加输入表示补曲率；这些表示可与不同损失或GLM结合，组合后的参数学习也需匹配。IRLS贯穿Huber回归与GLM：两者都在当前估计处把非线性损失线性化为加权最小二乘，区别只在权重来自稳健残差还是指数族方差；惩罚可以加到GLM得分上得到惩罚广义线性模型，核内积替换同时出现在SVR与KRR中，基函数预测器可接入GLM与GAM。组合时必须保持目标一致：先确定残差计分与响应族，再叠加约束与表示，不能把两种损失混为一谈，也不能把核空间惩罚误当作响应分布假设。

最后比较这些系数到底解释什么。OLS与GLM（恒等链接）中$\partial f/\partial x_j=w_j$，斜率就是该特征对预测的边际效应；对数链接GLM中$\mu=g^{-1}(\eta)$，边际效应为$\partial\mu/\partial x_j=\mu(1-\mu)\beta_j$（logit）或$\mu\beta_j$（log），系数$\beta_j$必须经链接反变换后才能解释成对响应的影响。PCR与PLS把系数先组合在子空间里：预测$f(\bm{x})=\widehat b+\bm{c}^{\top}\bm{V}_k^{\top}(\bm{x}-\bar{\bm{x}})$，边际效应为$\bm{V}_k\bm{c}$，单个原始特征的系数被分散到多个成分方向上。基函数回归中$w_j$作用在$\phi_j(\bm{x})$上，不是作用在原始输入上。模型名称相近不意味着系数含义相同，比较前先确认预测器经过了哪些变换与链接。
